\documentclass[12pt]{article}
\usepackage[margin=1in]{geometry}
\usepackage{enumitem}
\usepackage{titlesec}
\usepackage{parskip}
\usepackage{graphicx}
\usepackage{xcolor}
\usepackage{hyperref}
\usepackage{fontenc}

\title{\textbf{Acemoglu and Robinson 2001 RG Notes}\\
\large Acemoglu, Daron, and James A Robinson. ``A Theory of Political Transitions'' [2001].\\
\textit{American Economic Review} 91(4): 938-963.}
\author{}
\date{}

\begin{document}

\maketitle

General framework: The wealthy are faced with a threat of revolution/unrest and are the driving decision-makers in whether or not to allow for peaceful democratization or take the chances of violence.

Four primary conclusions from first model:

1) ``Societies with more initial inequality are more likely to switch between democracy and nondemocracy, and less likely to have a fully consolidated democracy. So our results are in line with the empirical findings of a positive association between inequality and political instability.''

2) ``There is a nonmonotonic relationship between inequality and redistribution, with societies at intermediate levels of inequality redistributing more than both very equal and very unequal societies.''

3) ``The relationship between fiscal volatility and inequality is likely to be increasing.''

4) ``In societies in which taxation creates less economic distortions --- for example, in societies where a large fraction of the GDP is generated from natural resources ---democracies may be harder to consolidate.''

\textit{\underline{``The role of recessions is to change the opportunity cost of coups to rich agents in a democracy and of revolution to poor agents in a nondemocracy.''}}

Democratization is more likely in periods of recession whereby the poor have little to lose in destroying the current assets/wealth in the case of revolution: if they are to receive very little anyways, then the eradication of that good becomes meaningless. Similarly, if inequality is sufficiently high, the poor have little incentive to avoid revolution even in times of non-recessions, as they still stand to lose very little.

\end{document}
